% Overleaf compiler: pdfLaTeX.
\RequirePackage{fix-cm}              % Scale the default Computer Modern fonts
\documentclass[a4paper]{article}
\usepackage[left=21mm,right=21mm,top=23mm,bottom=22mm]{geometry}
\usepackage{amsmath}                 % Aligned equations and named operators
\usepackage{amsfonts}                % Blackboard-bold symbols
\usepackage{multicol}                % Two columns below a full-width title
\usepackage{titlesec}                % Section heading style and spacing
\usepackage{enumitem}                % List margins and spacing
\usepackage[table]{xcolor}           % Colors and table support (loads array via colortbl)
\usepackage{booktabs}                % Table rules
\usepackage{listings}                % Python syntax highlighting
\usepackage[style=ieee]{biblatex}    % IEEE citations and bibliography
\usepackage[hidelinks]{hyperref}     % Linked equation and bibliography references

\addbibresource{references.bib}
\DeclareMathOperator*{\argmin}{arg\,min}

\definecolor{accent}{HTML}{194F54}
\definecolor{codebackground}{HTML}{F2F5F5}
\definecolor{codekeyword}{HTML}{D73948}
\definecolor{codefunction}{HTML}{4B69C6}
\definecolor{codenumber}{HTML}{B60157}

\pagestyle{empty}
\setlength{\columnsep}{10mm}
\setlength{\multicolsep}{20bp}
\titleformat{\section}{\normalfont\bfseries\fontsize{13bp}{16bp}\selectfont\color{accent}}{}{0pt}{}
\titlespacing*{\section}{0pt}{16bp}{4bp}
\setlist{leftmargin=13bp,labelwidth=8bp,labelsep=5bp,align=left,
  topsep=3bp,partopsep=0pt,itemsep=-1.5bp,parsep=0pt}
\setlength{\biblabelsep}{6bp}
\setlength{\bibitemsep}{3.5bp}
\AtBeginBibliography{\fontsize{9bp}{12.5bp}\selectfont}

\renewcommand{\footnotesize}{\fontsize{8.5bp}{11.5bp}\selectfont}
\renewcommand{\footnoterule}{%
  \kern-3bp\hrule width 32mm height 0.4bp\kern 2.6bp
}
\setlength{\skip\footins}{7bp}
\setlength{\footnotesep}{8.5bp}

\lstset{
  language=[3]Python,
  basicstyle=\fontsize{8.5bp}{11.5bp}\selectfont\ttfamily,
  keywordstyle=\color{codekeyword},
  emph={ridge,eye,solve,array,print},
  emphstyle=\color{codefunction},
  % Match Typst's colors for the characters used in this example.
  literate=
    {1}{{\textcolor{codenumber}{1}}}1
    {2}{{\textcolor{codenumber}{2}}}1
    {3}{{\textcolor{codenumber}{3}}}1
    {9}{{\textcolor{codenumber}{9}}}1
    {=}{{\textcolor{codekeyword}{=}}}1
    {@}{{\textcolor{codekeyword}{@}}}1
    {+}{{\textcolor{codekeyword}{+}}}1
    {*}{{\textcolor{codekeyword}{*}}}1
    {-}{{\textcolor{codekeyword}{-}}}1,
  columns=fullflexible,
  keepspaces=true,
  backgroundcolor=\color{codebackground},
  frame=single,
  framerule=0pt,
  framesep=7bp,
  xleftmargin=7bp,
  xrightmargin=7bp,
  aboveskip=7bp,
  belowskip=7bp
}

\begin{document}
\raggedright
\fontsize{10.5bp}{14bp}\selectfont
\setlength{\abovedisplayskip}{9bp}
\setlength{\belowdisplayskip}{9bp}
\setlength{\abovedisplayshortskip}{9bp}
\setlength{\belowdisplayshortskip}{9bp}

% Reserve the same title and subtitle height and depth as in ridge.typ.
\noindent\raisebox{0pt}[23.2bp][5.8bp]{%
  \fontsize{29bp}{29bp}\selectfont\bfseries Ridge regression}%
\par
\vspace{10bp}
\nointerlineskip
\noindent\raisebox{0pt}[9.6bp][2.4bp]{%
  \fontsize{12bp}{12bp}\selectfont\color{accent}%
  A worked example with two correlated predictors}%
\par
\vspace{16bp}
{\color{accent}\hrule height 0.7bp}

\raggedcolumns
\begin{multicols}{2}
\setlength{\parskip}{3.5bp}

\section*{The problem}

Two predictors can carry almost the same information. Ordinary least squares
can then assign them large coefficients of opposite signs. The predictions can
stay nearly the same even as the coefficients swing.

Ridge regression adds a penalty for large coefficients. Hoerl and Kennard
introduced the estimator~\cite{hoerl-estimator} and illustrated its behavior with
examples~\cite{hoerl-applications}.

\section*{The estimator}

Let \(X\) be the \(n\)-by-\(p\) matrix of predictors and \(y\) the response vector.
Take the columns and the response to be centered.\footnote{%
  Centering puts the intercept at zero, so the penalty applies only to the
  slopes.%
} For a penalty strength \(\lambda \ge 0\), the ridge estimate is

\begin{equation}
  \label{ridge}
  \begin{aligned}
    \hat{\beta}_\lambda
      &= \argmin_{\beta \in \mathbb{R}^p}
         \bigl[\lVert y - X\beta\rVert_2^2 + \lambda\lVert\beta\rVert_2^2\bigr] \\
      &= (X^\top X + \lambda I_p)^{-1} X^\top y.
  \end{aligned}
\end{equation}

For \(\lambda > 0\), the added diagonal makes the matrix invertible even if
two columns of \(X\) are identical. At \(\lambda = 0\), the inverse requires
full column rank.

\section*{Fitting the model}

\begin{enumerate}
  \item Center the predictors and the response.
  \item Pick a penalty strength \(\lambda\).
  \item Solve the system in Equation~\ref{ridge}.
\end{enumerate}

\section*{Reading the fit}

The first row of the table shows how a perfect fit can hide large slopes.
Consider the observation where both predictors are \(2\).

\begin{itemize}
  \item At \(\lambda = 0\), the slopes nearly cancel. The fitted value is
    \(2 (-15.50 + 17.00) = 3.00\).
  \item At \(\lambda = 1\), both slopes are below one. The fitted value drops to
    about \(2.26\).
\end{itemize}

\columnbreak

\section*{Computing the solution}

The code uses a linear solve to evaluate Equation~\ref{ridge} without forming a matrix
inverse. The two columns of \(X\) differ by just \(0.1\) in two of the four rows.

\begin{lstlisting}
import numpy as np

def ridge(X, y, lam):
    gram = X.T @ X
    identity = np.eye(X.shape[1])
    return np.linalg.solve(
        gram + lam * identity, X.T @ y
    )

X = np.array([[1, .9], [2, 2],
              [-1, -.9], [-2, -2]])
y = np.array([-.2, 3, .2, -3])

print(ridge(X, y, lam=1))
\end{lstlisting}

\section*{What changes}

The table reports both slopes and the root mean square error (RMSE) on the
same four observations used to fit the model.\footnote{%
  Training RMSE measures fit on these four rows. Choose the penalty using
  validation data when the goal is prediction.%
}

\begingroup
\setlength{\tabcolsep}{0pt}
\setlength{\aboverulesep}{0pt}
\setlength{\belowrulesep}{0pt}
\renewcommand{\arraystretch}{1.18}
\arrayrulecolor{accent}
\begin{tabular}{@{}p{.25\columnwidth}*{3}{>{\raggedleft\arraybackslash}p{.25\columnwidth}}@{}}
  \toprule[0.8bp]
  \(\lambda\) & \(\hat{\beta}_1\) & \(\hat{\beta}_2\) & \textbf{RMSE} \\
  \midrule[0.4bp]
  0.00 & \(-15.50\) & 17.00 & 0.00 \\
  0.01 &  \(-6.63\) &  7.96 & 0.57 \\
  0.10 &  \(-0.62\) &  1.82 & 0.96 \\
  1.00 &       0.44 &  0.69 & 1.03 \\
  \bottomrule[0.8bp]
\end{tabular}
\par
\endgroup

\section*{References}

\printbibliography[heading=none]
\end{multicols}
\end{document}
